% mage-auto-heading: problem
% mage-auto-heading: conceptual-design
% mage-auto-heading: engineering-models
% mage-auto-heading: realization
% mage-auto-heading: derive-the-selection
% mage-auto-heading: connect-selection-execution-and-interpretation
% mage-auto-heading: challenge-the-impact-declaration
% mage-auto-heading: introduce-the-arrangement-incrementally
% mage-auto-heading: adaptation-and-limits
% Evidence: book/data/local-evidence/blueprints-11202026/evidence-ledger.md; DocAble revision 563cdae2316df8798391d17c794ae7f662410321.
% Authoritative LaTeX manuscript; Markdown is a generated compatibility projection.

\section{Problem}\label{heading:problem}

As agents produce changes faster, the cost of evaluating those changes can become a bottleneck. Running every test and analysis after every change provides broad evidence, but much of that work may be unrelated to the modification. DocAble's full predeployment validation includes roughly four thousand tests and numerous static checks. Repeating all of them at every integration boundary would consume substantial computing resources and delay subsequent work.

Selecting fewer checks requires evidence that the reduced set is adequate. An omitted dependency can cause a relevant check to be skipped, allowing a regression to escape. Moreover, selecting checks is not the same as executing them: the validation machinery must actually honor the selection, and the acceptance decision must interpret the corresponding results.

\section{Conceptual Design}\label{heading:conceptual-design}

The engineering opportunity is to derive validation from models of change impact. A component model identifies the implementation elements involved; relationships among components identify potentially affected subsystems; declarations of validation coverage identify the checks to run. This arrangement can reduce repeated validation without requiring the engineer or agent to choose tests manually. The design below uses DocAble's change-impact selector as a concrete example.

\autoref{handbook-model-based-test-selection} illustrates the design. A proposed change is resolved through component and subsystem models to a declared set of potentially affected subsystems. Their validation obligations determine which tests and analyses should execute. Unmapped changes or recognized gaps in the impact model require broader validation. Separate analyses evaluate the selector itself, including comparisons with results from the full suite.

\begin{figure}[!htb]
\centering
\includegraphics[width=\linewidth]{book/latex/assets/handbook-model-based-test-selection-3cbe737b13.pdf}
\caption[Model-based regression-test selection]{\emph{Model-based regression-test selection.} Component and dependency models derive affected subsystems, and validation coverage maps them to checks. Execution records and findings feed an acceptance gate under an explicit policy. Dependency analyses and full-suite comparisons challenge the selection; conservative fallback covers recognized uncertainty, not every unmodeled dependency.}
\label{handbook-model-based-test-selection}
\end{figure}

% Evidence: BP-SELECTION | source: book/data/local-evidence/evidentiary-basis-11202026/integration/records/BP-SELECTION.md | kind: implementation-and-illustration
Suppose an agent changes the component responsible for editing PDF structure. The selector first resolves the modified files to their modeled components. It uses their declared subsystem memberships and any rules that broaden the affected set to identify potentially affected subsystems, then consults the validation-coverage model to determine which checks evaluate those subsystems. For a change whose modeled impact is confined to PDF processing, PDF-specific checks may be selected while unrelated spreadsheet checks are omitted. The executor must honor that selection, and the acceptance gate must interpret the results of the checks actually required.

The decision is only as reliable as the relationships behind it. If the changed component also affects a shared document abstraction, the selection must cover the affected consumers, whether through a widened declaration or a broader run. An unmapped dependency must not be interpreted as proof that no additional checks are needed. Separately recovered dependency relationships and comparisons with full-suite execution can expose omissions in the selection model. The blueprint therefore governs not only which checks run, but how engineers establish confidence that the selection remains adequate.

\section{Engineering Models}\label{heading:engineering-models}

Five models connect a change to validation. Their outputs have different types: changed files identify implementation elements; impact analysis declares subsystems; coverage declarations map subsystems to validation units. Keeping these distinctions explicit prevents a plausible dependency analysis from being mistaken for an adequate test selection.

The model relationships are summarized in \autoref{handbook-selection-models}.

{\footnotesize
\begin{longtable}[]{@{}
  >{\raggedright\arraybackslash}p{(\linewidth - 6\tabcolsep) * \real{0.2500}}
  >{\raggedright\arraybackslash}p{(\linewidth - 6\tabcolsep) * \real{0.2500}}
  >{\raggedright\arraybackslash}p{(\linewidth - 6\tabcolsep) * \real{0.2500}}
  >{\raggedright\arraybackslash}p{(\linewidth - 6\tabcolsep) * \real{0.2500}}@{}}
\caption[Engineering models for regression-test selection]{\emph{Engineering models for regression-test selection.} Component and dependency relationships derive an affected-subsystem declaration; validation coverage maps it to checks. The execution model connects that selection to actual work and acceptance, while separate evidence challenges omissions.}\label{handbook-selection-models}\tabularnewline
\toprule\noalign{}
\begin{minipage}[b]{\linewidth}\raggedright
Model
\end{minipage} & \begin{minipage}[b]{\linewidth}\raggedright
Engineering question
\end{minipage} & \begin{minipage}[b]{\linewidth}\raggedright
Representation
\end{minipage} & \begin{minipage}[b]{\linewidth}\raggedright
Relationships / consumers
\end{minipage} \\
\midrule\noalign{}
\endfirsthead
\toprule\noalign{}
\begin{minipage}[b]{\linewidth}\raggedright
Model
\end{minipage} & \begin{minipage}[b]{\linewidth}\raggedright
Engineering question
\end{minipage} & \begin{minipage}[b]{\linewidth}\raggedright
Representation
\end{minipage} & \begin{minipage}[b]{\linewidth}\raggedright
Relationships / consumers
\end{minipage} \\
\midrule\noalign{}
\endhead
\bottomrule\noalign{}
\endlastfoot
Component model & Where do changed implementation elements belong? & File-to-component identities and component-to-subsystem membership. & The impact query reuses the component resolver; stable identities join the other models. \\
Dependency model & Where could a change have consequences? & Declared component relationships and separately recovered project/import dependencies. & Declared relationships inform impact; recovered edges challenge the declaration. \\
Validation coverage model & Which checks evaluate each obligation or subsystem? & Validation-unit scope labels and supporting links to obligations or coverage observations. & Unit labels translate affected subsystems into checks; observations challenge coverage claims. \\
Change-impact model & Which subsystems may this change affect? & Derived affected-set declaration, derivation paths, widening hints, full-run requirement. & Computed from changed files and maintained relationships; consumed by unit selection and selector validation. \\
Validation execution model & What must run, and how do findings govern acceptance? & Selected check identities, execution modes, required results, and acceptance policy. & Executor consumes the selection; execution records support correspondence checks and gate decisions. \\
\end{longtable}
}

% Evidence: BP-SELECTION | source: book/data/local-evidence/evidentiary-basis-11202026/integration/records/BP-SELECTION.md | kind: implementation-and-design
A change-impact model relates implementation elements to the subsystems and obligations they may affect. DocAble's component model identifies ownership and structural relationships within its codebase. A validation model groups components into subsystems, including PDF processing, Office document processing, the web application, and deployment infrastructure. The selector reuses the component model's path resolver rather than maintaining a second account of component membership.

% Evidence: BP-SELECTION | source: book/data/local-evidence/evidentiary-basis-11202026/integration/records/BP-SELECTION.md | kind: implementation-and-design
Given a set of changed files, the selector resolves each file to its component and derives the declared affected subsystems. Each validation unit declares the subsystems it evaluates. The selection rule includes a unit when its scope intersects the affected set. This is a decision about what should execute; the executor must consume it.

The \textbf{dependency model} has two uses that must remain distinguishable. Maintained relationships express the impact assumptions used for selection. Separately recovered project and import relationships provide descriptive evidence that can challenge those assumptions. Both may share component identities, so agreement is useful without constituting independent proof of completeness.

The \textbf{change-impact model} is a derived view for a particular change: an affected-subsystem declaration together with its derivation and any requirement for full validation. Derivation does not by itself confer authority. The engineer decides where the evidence justifies allowing that declaration to control a reduced run. The \textbf{validation execution model} then specifies required execution and the policy for interpreting findings. Execution records describe what actually executes; compare them with the selected set and the permitted broader modes.

Give component ownership, dependency declarations, and validation-unit scope explicit maintainers. A file move must preserve or update its component identity; a new cross-component call may widen impact; a check that acquires another responsibility needs a revised scope declaration. Observed coverage can challenge those declarations but cannot establish that an unobserved relationship is absent. These correspondence obligations keep the joins useful as the system changes.

The executor can choose scheduling and parallelism where the declared dependencies permit. It may run a deliberately broader set under the acceptance policy. Neither choice authorizes silently dropping a required check or changing what a finding means for admission.

\section{Realization}\label{heading:realization}

\subsection*{Derive the Selection}\label{heading:derive-the-selection}

The following explanatory pseudocode preserves DocAble's selection decision order, with type annotations omitted and diagnostic messages simplified:

% mage-code-language: python
\begin{verbatim}
def decide_unit(unit_labels, affected, forced_full):
    if forced_full:
        return True, "full validation required"
    if not unit_labels:
        return True, "validation unit has no scope declaration"
    if unit_labels & affected:
        return True, "unit covers an affected subsystem"
    return False, "no modeled impact"
\end{verbatim}

The unconditional branches retain checks when the change requires full validation or the validation unit lacks a scope declaration.

% Evidence: BP-SELECTION | source: book/data/local-evidence/evidentiary-basis-11202026/integration/records/BP-SELECTION.md | kind: implementation-and-design
DocAble's selector returns a full-validation requirement for unmapped or cross-cutting files and dependency-manifest changes. The latter deserve special treatment because a project dependency can affect the assembled product beyond the component containing the manifest. Optional impact hints can enlarge the declared affected set, but cannot remove checks required by the models. These fallbacks protect against recognized uncertainty; an unrecognized dependency can still cause an omission.

\subsection*{Connect Selection, Execution, and Interpretation}\label{heading:connect-selection-execution-and-interpretation}

% Evidence: E59 | source: book/data/local-evidence/evidentiary-basis-11202026/integration/records/E59.md | kind: dated-case-evidence
A DocAble merge check illustrates a failure to connect selection to execution. This separate selector chooses inexpensive lints for merge validation. In the original configuration, a full-tree walker still executes hundreds of predicates regardless of that selection, accounting for almost all validation time in a profiled run.

% Evidence: E59 | source: book/data/local-evidence/evidentiary-basis-11202026/integration/records/E59.md | kind: dated-case-evidence
Findings outside the selected set can also affect the gate's verdict. Selection therefore governs neither the work performed nor all the evidence interpreted.

% Evidence: E59 | source: book/data/local-evidence/evidentiary-basis-11202026/integration/records/E59.md | kind: dated-case-evidence
The repair connects the walker to the same selection as the other lint processes, with three cases for its own cohort of checks:

\begin{itemize}
\tightlist
\item
  \textbf{No restriction:} execute the ordinary full walker.
\item
  \textbf{Empty intersection:} no walker checks are selected, so skip the walker.
\item
  \textbf{Nonempty intersection:} at least one walker check is selected, so execute the full predicate cohort.
\end{itemize}

% Evidence: E59 | source: book/data/local-evidence/evidentiary-basis-11202026/integration/records/E59.md | kind: dated-case-evidence
The nonempty-intersection case deliberately retains the full cohort to avoid omitting a selected predicate. The empty-intersection branch removes both the unnecessary walk and its contribution to the gate verdict.

% Evidence: E59 | source: book/data/local-evidence/evidentiary-basis-11202026/integration/records/E59.md | kind: dated-case-evidence
Historical replay shows a substantial reduction in validation time for the evaluated merge workload: the revised configuration completes the batches while the earlier configuration reaches the execution limit. This comparison does not establish a general performance rate.

The incident distinguishes \textbf{selection}, \textbf{execution}, and \textbf{interpretation}. Selection names required checks; execution performs them or an explicitly broader set; interpretation determines how findings affect acceptance. Unexpected findings require the declared blocking, escalation, or review response. Test these responsibilities together, including cases where broader execution produces an unexpected failure.

\subsection*{Challenge the Impact Declaration}\label{heading:challenge-the-impact-declaration}

The selection must supply the required evidence for every validation obligation the change can affect. Because the impact model can omit relationships, DocAble's validation machinery challenges its declarations with several sources of evidence.

% Evidence: BP-SELECTION | source: book/data/local-evidence/evidentiary-basis-11202026/integration/records/BP-SELECTION.md | kind: implementation-and-design
Static dependency analysis constructs a separate estimate of potentially affected subsystems. For C\# code, it follows reverse project references; for Python, it follows reverse module imports. The validation condition requires the subsystems reached by this analysis to be included in the selector's declaration:

\[
\operatorname{Reachable}(C) \subseteq \operatorname{Declared}(C).
\]

Both sets contain \textbf{validation subsystems} for a proposed change C. An omitted reachable subsystem is a counterexample to the impact claim; validation-unit scope declarations then translate subsystems into checks.

The reachable set under-approximates potential impact: the graphs may omit reflection, dynamic dispatch, or configuration-dependent behavior. Passing the subset check establishes consistency with the dependencies the analysis finds. Both calculations also share the component model, so a common classification error can escape the comparison.

% Evidence: BP-SELECTION | source: book/data/local-evidence/evidentiary-basis-11202026/integration/records/BP-SELECTION.md | kind: implementation-and-design
Additional model relationships provide another source of counterexamples. DocAble's validator compares the declared subsystems with maintained relationships among controls, mutations, and other engineering mechanisms. These checks can expose dependencies not visible in ordinary project-reference or import graphs.

Coverage evidence supplies another challenge. If an execution associates a validation unit with affected implementation, but the selector omits that unit, the disagreement warrants investigation. An unobserved path cannot establish the absence of a dependency.

Full-suite execution supplies an empirical comparison. Record which checks the selector omits, then investigate failures among them. Determine whether the change causes the failure, whether the check's scope is accurate, and whether the result is reproducible; pre-existing or intermittent failures need not imply a new dependency.

% Evidence: BP-SELECTION | source: book/data/local-evidence/evidentiary-basis-11202026/integration/records/BP-SELECTION.md | kind: implementation-and-design
Historical evaluations expose under-declared dependencies that widening the impact model resolves on the evaluated corpus. A subsequent evaluation over a larger, different collection of changes finds additional omissions. Eliminating findings on one corpus therefore does not establish adequacy for future changes. Repeated evaluation seeks counterexamples to declared impact relationships rather than certifying their completeness.

A discrepancy calls for revising the impact model and repeating the comparison. Widening the selected set cannot introduce another omission for the same change, although it increases validation cost.

\subsection*{Introduce the Arrangement Incrementally}\label{heading:introduce-the-arrangement-incrementally}

Begin with a full validation process whose results are understood. Partition the system into components and validation subsystems, then associate checks with each subsystem. Coarse partitions run more checks but require fewer assumptions about independence.

Implement the selector as a query that explains each file's component, the relationships establishing affected subsystems, and why each check is included or omitted. These records make mistaken selections diagnosable.

Verify selection, execution, and gate decisions together. Exercise empty and full selections and checks invoked through different mechanisms. Inject failures inside and outside the selected set, including missing results, and confirm the intended blocking, broader-run, or review response. Compare all records against the same revision.

Keep selection advisory while evaluating historical changes and current work alongside full validation. Resolve omissions, then repeat the comparison on varied workloads, including shared dependencies, test infrastructure, and configuration.

Then decide whether the evidence justifies reduced validation at a particular boundary. A team may use selected checks for early integration while retaining full validation before release and periodic full-suite comparisons. The decision depends on the consequences of an escaped regression and the properties the baseline checks actually evaluate.

Maintainers must revise the models and repeat these comparisons as component boundaries, test responsibilities, and dependencies change.

\section{Adaptation and Limits}\label{heading:adaptation-and-limits}

Reduced repeated validation must justify the costs of maintaining impact relationships, evaluating the selector, investigating omissions, and running broader checks. The lifecycle account in \autoref{6.1-toward-a-theory-of-mage:heading:cost-and-failure-separately} supports that comparison while preserving the required quality of acceptance evidence.

The comparison depends on actual work executed, including validation latency, computing cost, omitted-check findings, and model maintenance. A small selected set offers little saving if its executor still runs broad checks; a cheap run can also defer detection or permit a defect to escape.

The appropriate granularity depends on the system. A small project may gain little from selective validation because its complete test suite is inexpensive. A large repository with costly integration tests may justify a detailed model of change impact. Where the consequences of an omission are severe, conservative selection or full validation may remain preferable.

The same method can select performance, security, or deployment checks when models connect proposed changes to the obligations those checks evaluate.
